\subsection{Linear Map Operations and Examples}
% **Role:** algebraic language of linear structure.
%
% Linear maps, kernels, ranges, rank-nullity, isomorphisms, matrices of maps, change of basis, similarity.

Linear maps preserve vector-space structure. Kernels, ranges, rank-nullity, matrices, change of basis, and similarity describe them without making coordinates primary.

\begin{definition} \label{def:linear-map} \hlt{Linear Map}\\
Let $\mathcal{V}$ and $\mathcal{W}$ be vector spaces over the same field $\mathbb{F}$.\\
A function $T:\mathcal{V}\to\mathcal{W}$ is \hlt{linear} if $T(a\mathbf{u}+b\mathbf{v})=aT\mathbf{u}+bT\mathbf{v}$ for all $\mathbf{u},\mathbf{v}\in\mathcal{V}$ and $a,b\in\mathbb{F}$.\\
The set of all linear maps from $\mathcal{V}$ to $\mathcal{W}$ is denoted $\LL(\mathcal{V},\mathcal{W})$. If $\mathcal{V}=\mathcal{W}$, elements of $\LL(\mathcal{V})$ are \hlt{operators} on $\mathcal{V}$.
\end{definition}

\begin{definition} \hlt{Zero and Identity Maps}\\
For vector spaces $\mathcal{V}$ and $\mathcal{W}$, the \hlt{zero map} from $\mathcal{V}$ to $\mathcal{W}$ is $0_{\mathcal{V},\mathcal{W}}:\mathcal{V}\to\mathcal{W}$, where $0_{\mathcal{V},\mathcal{W}}\mathbf{v}=\mathbf{0}_{\mathcal{W}}$ for all $\mathbf{v}\in\mathcal{V}$.\\
The \hlt{identity map} on $\mathcal{V}$ is $I_{\mathcal{V}}:\mathcal{V}\to\mathcal{V}$, where $I_{\mathcal{V}}\mathbf{v}=\mathbf{v}$ for all $\mathbf{v}\in\mathcal{V}$.
\end{definition}

\begin{proposition} \label{prop:linear-maps-send-zero-to-zero} \hlt{Linear Maps Send Zero to Zero}\\
If $T\in\LL(\mathcal{V},\mathcal{W})$, then $T\mathbf{0}_{\mathcal{V}}=\mathbf{0}_{\mathcal{W}}$.
\begin{proof}
Linearity gives $T\mathbf{0}_{\mathcal{V}}=T(\mathbf{0}_{\mathcal{V}}+\mathbf{0}_{\mathcal{V}})=T\mathbf{0}_{\mathcal{V}}+T\mathbf{0}_{\mathcal{V}}$. Adding $-T\mathbf{0}_{\mathcal{V}}$ to both sides gives $T\mathbf{0}_{\mathcal{V}}=\mathbf{0}_{\mathcal{W}}$.\\
\hlt{Structural Commentary}: Linearity preserves zero.\\
\hlt{Key Insight}: A linear map cannot move zero away from zero.
\end{proof}
\end{proposition}

\begin{definition} \hlt{Operations on Linear Maps}\\
For $S,T\in\LL(\mathcal{V},\mathcal{W})$ and $a\in\mathbb{F}$, define pointwise $(S+T)\mathbf{v}\defeq S\mathbf{v}+T\mathbf{v}$ and $(aT)\mathbf{v}\defeq a(T\mathbf{v})$ for every $\mathbf{v}\in\mathcal{V}$.
\end{definition}

\begin{proposition} \label{prop:linear-maps-form-vector-space} \hlt{Linear Maps Form a Vector Space}\\
With the operations above, $\LL(\mathcal{V},\mathcal{W})$ is a vector space over $\mathbb{F}$.
\begin{proof}
If $S,T\in\LL(\mathcal{V},\mathcal{W})$ and $a,c\in\mathbb{F}$, then for $\mathbf{u},\mathbf{v}\in\mathcal{V}$,
\begin{align}
(aS+T)(c\mathbf{u}+\mathbf{v})
=a(cS\mathbf{u}+S\mathbf{v})+(cT\mathbf{u}+T\mathbf{v})
=c(aS+T)\mathbf{u}+(aS+T)\mathbf{v}. \nonumber
\end{align}
Thus $\LL(\mathcal{V},\mathcal{W})$ is closed under linear combinations. The zero map is the additive identity, $-T$ is the additive inverse of $T$, and the remaining vector-space axioms hold pointwise because they hold in $\mathcal{W}$.\\
\hlt{Structural Commentary}: Linear map operations are pointwise.\\
\hlt{Key Insight}: Pointwise sums and scalar multiples preserve linearity.
\end{proof}
\end{proposition}

\begin{definition} \hlt{Product of Linear Maps}\\
For $T\in\LL(\mathcal{U},\mathcal{V})$ and $S\in\LL(\mathcal{V},\mathcal{W})$, the \hlt{product} $ST\in\LL(\mathcal{U},\mathcal{W})$ is defined by $(ST)\mathbf{u}\defeq S(T\mathbf{u})$ for all $\mathbf{u}\in\mathcal{U}$.
\end{definition}

\begin{proposition} \label{prop:linear-map-product-algebra} \hlt{Algebra of Linear Map Products}\\
If $T\in\LL(\mathcal{U},\mathcal{V})$ and $S\in\LL(\mathcal{V},\mathcal{W})$, then $ST\in\LL(\mathcal{U},\mathcal{W})$. Whenever defined, $(RS)T=R(ST)$, $TI_{\mathcal{U}}=T=I_{\mathcal{V}}T$, $(S_1+S_2)T=S_1T+S_2T$, and $S(T_1+T_2)=ST_1+ST_2$.
\begin{proof}
For $\mathbf{u},\mathbf{x}\in\mathcal{U}$ and $a\in\mathbb{F}$, $(ST)(a\mathbf{u}+\mathbf{x})=S(aT\mathbf{u}+T\mathbf{x})=aS(T\mathbf{u})+S(T\mathbf{x})=a(ST)\mathbf{u}+(ST)\mathbf{x}$, hence $ST\in\LL(\mathcal{U},\mathcal{W})$. For each remaining identity $A=B$ and admissible $\mathbf{z}$, $A\mathbf{z}=B\mathbf{z}$ by associativity of evaluation, identity maps, or distributivity in the target space; hence $A=B$.\\
\hlt{Structural Commentary}: Composition chains linear actions, while addition distributes through evaluation.\\
\hlt{Key Insight}: Linear map algebra is function composition plus pointwise vector-space algebra.
\end{proof}
\end{proposition}

\begin{definition} \hlt{Graph of a Function}\\
For a function $T:\mathcal{V}\to\mathcal{W}$, the \hlt{graph} of $T$ is $\operatorname{graph}(T)\defeq\{(\mathbf{v},T\mathbf{v}):\mathbf{v}\in\mathcal{V}\}\subseteq\mathcal{V}\times\mathcal{W}$.
\end{definition}

\begin{proposition} \label{prop:graph-subspace-characterizes-linearity} \hlt{Linearity and Graphs}\\
A function $T:\mathcal{V}\to\mathcal{W}$ is linear if and only if $\operatorname{graph}(T)$ is a subspace of $\mathcal{V}\times\mathcal{W}$.
\begin{proof}
\begin{enumerate}[label=(\roman*)]
\setlength{\itemsep}{0pt}
\item[$(\Rightarrow)$] If $T$ is linear, then $(\mathbf{0},\mathbf{0})\in\operatorname{graph}(T)$.\\
If $(\mathbf{u},T\mathbf{u}),(\mathbf{v},T\mathbf{v})\in\operatorname{graph}(T)$ and $a\in\mathbb{F}$, then $a(\mathbf{u},T\mathbf{u})+(\mathbf{v},T\mathbf{v})=(a\mathbf{u}+\mathbf{v},T(a\mathbf{u}+\mathbf{v}))\in\operatorname{graph}(T)$.\\
Thus $\operatorname{graph}(T)$ is a subspace.
\item[$(\Leftarrow)$] If $\operatorname{graph}(T)$ is a subspace, then for $\mathbf{u},\mathbf{v}\in\mathcal{V}$ and $a\in\mathbb{F}$, the vector
$a(\mathbf{u},T\mathbf{u})+(\mathbf{v},T\mathbf{v})=(a\mathbf{u}+\mathbf{v},aT\mathbf{u}+T\mathbf{v})$
lies in $\operatorname{graph}(T)$. Hence $T(a\mathbf{u}+\mathbf{v})=aT\mathbf{u}+T\mathbf{v}$, so $T$ is linear.
\end{enumerate}
\hlt{Structural Commentary}: A graph records a function as a subset of a product space.\\
\hlt{Key Insight}: Closure of the graph is exactly linearity of the function.
\end{proof}
\end{proposition}

\begin{principle} \hlt{Coordinate-Free Principle}\\
A linear map is completely determined by what it does to a basis. Coordinates do not create the map; they only encode it once bases have been chosen.
\end{principle}

\begin{proposition} \label{prop:linear-map-determined-by-basis} \hlt{Linear Maps Are Determined by a Basis}\\
If $(\mathbf{v}_1, \ldots, \mathbf{v}_n)$ is a basis of $\mathcal{V}$ and $\mathbf{w}_1, \ldots, \mathbf{w}_n \in \mathcal{W}$ are arbitrary, then $\exists$ a unique linear map $T : \mathcal{V} \to \mathcal{W}$ such that $T(\mathbf{v}_j) = \mathbf{w}_j$ for $j=1,\ldots,n$.
\begin{proof}
Every $\mathbf{v} \in \mathcal{V}$ has a unique expression $\mathbf{v} = a_1\mathbf{v}_1 + \cdots + a_n\mathbf{v}_n$. Define $T(\mathbf{v}) \defeq a_1\mathbf{w}_1 + \cdots + a_n\mathbf{w}_n$.\\
This is well-defined because basis coordinates are unique, and linearity follows from coordinate expansions. Any linear map with the prescribed basis values must agree with this formula, so it is unique.\\
\hlt{Structural Commentary}: A basis reduces a linear map to finitely many choices.\\
\hlt{Key Insight}: Once basis images are fixed, linearity determines every other value.
\end{proof}
\end{proposition}

